\documentclass{article}
\usepackage{amsmath, amssymb, amsthm}
\usepackage{hyperref}
\usepackage{geometry}
\geometry{margin=1in}

\title{Erd\H{o}s Problem \#91}
\date{Last updated: 2025-08-31}
\author{Source: erdosproblems.com}

\begin{document}
\maketitle

\noindent\textbf{Status:} OPEN \quad \textbf{No prize}

\noindent\textbf{Tags:} geometry, distances

\medskip

\noindent\textbf{Problem Statement:}

Let $n$ be a sufficiently large integer. Suppose $A\subset \mathbb{R}^2$ has $\lvert A\rvert=n$ and minimises the number of distinct distances between points in $A$. Prove that there are at least two (and probably many) such $A$ which are non-similar.




For $n=3$ the equilateral triangle is the only such set. For $n=4$ the square or two equilateral triangles sharing an edge give two non-similar examples. For $n=5$ the regular pentagon is the unique such set (which has two distinct distances). Erd\H{o}s mysteriously remarks in \cite{Er90} this was proved by 'a colleague'. (In \cite{Er87b} this is described as 'a colleague from Zagreb (unfortunately I do not have his letter)'.) A published proof of this fact is provided by Kov\'{a}cs \cite{Ko24c}. In \cite{Er87b} Erd\H{o}s says that there are at least two non-similar examples for $6\leq n\leq 9$. The minimal possible number of distinct distances is the subject of [89] .




References

[Er87b] Erd\H{o}s, P., Some combinatorial and metric problems in geometry . Intuitive geometry (Si\'{o}fok, 1985) (1987), 167-177.

[Er90] Erd\H{o}s, Paul, Some of my favourite unsolved problems . A tribute to Paul Erd\H{o}s (1990), 467-478.

[Ko24c] Z. Kov\'{a}cs, A note on Erd\H{o}s's mysterious remark . arXiv:2412.05190 (2024).

\noindent\textbf{Related OEIS sequences:} A186704, possible


\bigskip
\noindent\small{Source: \url{https://www.erdosproblems.com/91}}
\end{document}
